\documentclass[10pt,a4paper]{amsart}
\usepackage[margin=19mm]{geometry}
\usepackage{amsmath,amssymb,mathtools}
\usepackage[scr=rsfs,cal=euler]{mathalfa}
\usepackage{xfrac}
\usepackage[unicode,hidelinks,bookmarks=false]{hyperref}
\newcommand{\ie}{i.e.\ }
\newcommand{\cf}{cf.\ }
\newcommand{\resp}{respectively\ }
\newcommand{\Subsection}{\subsection}
\providecommand{\nobreakdash}{\nobreak\hbox{-}}
\setlength{\emergencystretch}{2em}
\multlinegap0pt
\usepackage[shortlabels]{enumitem}
\usepackage{amssymb}
\usepackage{mathtools}
\usepackage{float}
\usepackage{tikz}
\newtheorem{definition}{Definition}
\newtheorem{lemma}{Lemma}
\newcommand{\Dp}{\mathcal{D}_P}
\newcommand{\F}{\mathcal{F}}
\newcommand{\N}{\mathbb{N}}
\newcommand{\e}{\epsilon}
  \def\emptyset{}%
\title{On minimal collections of sequences for testing continuity}
\author{Gyuhyun Lim}
\date{Source: arXiv:2605.11397v1 (CC BY 4.0)}
\begin{document}
\maketitle

\noindent\emph{Source and licence.} On minimal collections of sequences for testing continuity. Version arXiv:2605.11397v1 (CC BY 4.0), distributed by arXiv under the Creative Commons Attribution 4.0 International licence, http://creativecommons.org/licenses/by/4.0/, as recorded in the arXiv licence metadata for this exact version and shown on the abstract page https://arxiv.org/abs/2605.11397v1. Full licence notice: \texttt{licenses/arxiv-2605.11397-CC-BY-4.0.txt}.\par\smallskip

\noindent\textit{Editorial scope.} Counted unit: the article's lemma lem:Amin\_is\_testset (source0.tex lines 995--998). The article's complete original proof of it is delivered with it (source0.tex lines 1000--1048, one \texttt{proof} environment of 49 lines). No mathematics is written, repaired or re-derived: the statement and the proof are the article's own lines, in the article's own notation, and no retained line is substituted or re-flowed.  Retained uncounted as the source-local context the proof needs: lines 809--809; lines 929--938.  Nothing was omitted from any retained span, so the package carries no omission record.  No retained line contains a citation, so the wrapper carries no bibliography and no bibliography entry.  No retained line contains a figure environment, an image-inclusion command or drawing code, so no image is delivered and the package declares no asset dependency.  Editorial formatting only, in addition: an AMS \texttt{amsart} preamble that restates the article's own declarations and macros used by the retained lines, namely the environments \texttt{definition}, \texttt{lemma} on separate counters, together with the standard packages those lines need.  The retained environments print in this document as Definition 1, Lemma 1 in order of appearance, whereas the article prints the same items with the numbering of its own full document.  The complete native source of the article as distributed in the arXiv e-print for this exact version is bundled unchanged as \texttt{sources/200355/source0.tex}.
\section{Constructing a minimal test set}\label{sec:minimal}

\begin{definition}[Maximal almost disjoint family]\label{def:MAD}
A family $\mathcal M\subseteq\mathcal I_P$ is called a \emph{maximal almost disjoint family}, or simply a \emph{MAD family}, if
\begin{enumerate}[label=(\roman*)]
\item for $M_1\neq M_2$ in $\mathcal M$, one has $M_1\perp_P M_2$;
\item for every $M\in\mathcal I_P$, there exists $N\in\mathcal M$ such that
\[
|M\cap N|=\infty.
\]
\end{enumerate}
\end{definition}

\begin{lemma}\label{lem:Amin_is_testset}
Assume the standing framework of Section~\ref{sec:minimal}.
Then $A_{\min}$ is a test set, hence $A_{\min}\in\F$.
\end{lemma}

\begin{proof}
Let $f\in\Dp$.
Choose a witnessing sequence
\[
S=(S_n)\in D(f).
\]
Since $S\in D(f)$, there exists $\e>0$ such that the set
\[
I:=\{n\in\N:\ |f(S_n)-f(P)|\ge\e\}
\]
is infinite.
Choose strictly increasing indices $(n_k)_{k\in\N}\subseteq I$ and set
\[
R_k:=S_{n_k}.
\]
Let
\[
R:=\{R_k:k\in\N\}.
\]
We claim that $R\in\mathcal I_P$.

Since $R$ is the image of $\N$, it is at most countable.
If $R$ were finite, then $(R_k)$ would take values in a finite set $F$.
For each $x\in F$, choose an open neighborhood $U_x\in\mathcal N(P)$ with $x\notin U_x$.
Then
\[
U:=\bigcap_{x\in F}U_x
\]
is a neighborhood of $P$ disjoint from $F$.
But $(R_k)\to P$, so eventually $R_k\in U$, a contradiction.
Hence $R$ is infinite and therefore countably infinite.

Let $U\in\mathcal N(P)$.
Since $(R_k)\to P$, all but finitely many $R_k$ lie in $U$, so $R\setminus U$ is finite.
Thus $R\in\mathcal I_P$.

By Definition~\ref{def:MAD}(ii), choose $M\in\mathcal M$ with $|R\cap M|=\infty$.
Let $T^M\in A_{\min}$ enumerate $M$ injectively.
Since $R\cap M$ is infinite and $T^M$ visits each point of $M$ exactly once, the sequence $T^M$ passes through points of $R$ infinitely many times.
For those terms we have
\[
|f(T^M_j)-f(P)|\ge \e.
\]
Hence $f(T^M_n)\not\to f(P)$, so $T^M\in D(f)$.
Thus $A_{\min}\cap D(f)\neq\emptyset$.

Since $f\in\Dp$ was arbitrary, $A_{\min}$ is a test set.
Therefore $A_{\min}\in\F$.
\end{proof}

\end{document}
